% META: {"id":"1SPE-GEOREP-CO-003","chapitre":"1SPE-GEOMETRIE-REPEREE","type_objet":"corrige","exercice_ref":"1SPE-GEOREP-EX-003","capacites_codes":["C1"],"sources_inspiration":[],"status":"approved","origin":"generated","status_history":["generated","audited","accepted","approved"],"release_acceptance":"RELEASE_OWNER_BATCH_ACCEPTANCE","acceptance_closure_digest":"sha256:9b3ccf9a81c5520fbb7e03b7b2d3e7bf2057a02834b6d908bf3be5f49f3b553f","acceptance_receipt_digest":"sha256:4fad86f0282e0fa4e0419cca1ad6379ae7af5a280c04a4a90880f90a1c044242"}
% BEGIN-VERIFY
% from sympy import *
% x, y = symbols('x y')
% # Droite par A(0,4) de vecteur directeur u(3,2)
% # Vecteur normal n(2,-3), equation: 2(x-0) - 3(y-4) = 0 => 2x - 3y + 12 = 0
% assert 2*0 - 3*4 + 12 == 0
% # Verification avec un autre point: x=3, y = (2*3+12)/3 = 6
% assert 2*3 - 3*6 + 12 == 0
% END-VERIFY
\begin{corrige}{1SPE-GEOREP-EX-003}

\commentaireMarge{Si $\vec{u}(3 ; 2)$ est directeur, alors $\vec{n}(2 ; -3)$ est normal.}

\textbf{} Le vecteur directeur est $\vec{u}\begin{pmatrix} 3 \\ 2 \end{pmatrix}$, donc un vecteur normal est $\vec{n}\begin{pmatrix} 2 \\ -3 \end{pmatrix}$.

L'équation est $2(x - 0) - 3(y - 4) = 0$, soit $\boxed{2x - 3y + 12 = 0}$.

Vérification : $A(0 ; 4)$ : $2 \times 0 - 3 \times 4 + 12 = 0$ \checkmark.

\end{corrige}
